% Generated draft: requires mathematical and source review.
\documentclass[11pt]{article}
\usepackage[margin=25mm]{geometry}
\usepackage{fontspec}
\setmainfont{Latin Modern Roman}
\usepackage{amsmath,amssymb,mathtools,unicode-math}
\setmathfont{Latin Modern Math}
\usepackage{enumitem,needspace,xurl,hyperref}
\setlength{\parindent}{0pt}
\setlength{\parskip}{5pt}
\setlength{\emergencystretch}{3em}
\newcommand{\N}{\mathbb{N}}
\newcommand{\Z}{\mathbb{Z}}
\newcommand{\Q}{\mathbb{Q}}
\newcommand{\R}{\mathbb{R}}
\newcommand{\C}{\mathbb{C}}
\newcommand{\F}{\mathbb{F}}
\newcommand{\A}{\mathbb{A}}
\newcommand{\PP}{\mathbb{P}}
\newcommand{\E}{\mathbb{E}}
\newcommand{\eps}{\varepsilon}
\newcommand{\dd}{\,\mathrm{d}}
\begin{document}
\section*{UnsolvedMath ID: 11000033}
\textbf{Source problem code:} catalogue entry
\subsection{Definitions and mathematical statement}
A group \(H\) is freely indecomposable if it admits no decomposition \(H\cong A*B\) with both factors nontrivial. A nontrivial splitting over \(\Z\), in the usual Bass--Serre convention, is an amalgamated-product or HNN splitting with infinite cyclic edge group and nontrivial associated tree action, so a presentation with a collapsible edge is not counted merely as a formal splitting.

\emph{Question 3.12.} Is there an algorithm deciding, for a given subgroup \(H<\operatorname{Mod}_g\), whether \(H\) is freely indecomposable? Is there an algorithm deciding whether \(H\) splits over \(\Z\)? These are two decision questions. In a finite-generator input model, the input would be a list of words in fixed ambient generators whose subgroup is \(H\), and each procedure must halt with the correct answer on every valid input in its promised class. A finite-presentation input model or additional subgroup hypotheses could lead to a different problem.

The supplied statement does not specify such an effective encoding, assume finite generation or finite presentability, fix the genus and surface convention, or define the exact allowed cyclic splitting. The finite-generator description above is therefore a possible precise input model, not a silently added hypothesis. The question asks for decision procedures, not just enumeration of candidate splittings or a classification up to abstract isomorphism.
\subsection{Short English statement}
Decide algorithmically whether a specified mapping-class-group subgroup has a nontrivial free-product decomposition, and whether it splits over an infinite cyclic group. The source omits how the subgroup is given and any finiteness promises, which are essential to the algorithmic formulation.
\subsection{Sources}
\begin{itemize}
\item Local UnsolvedMath registry, permanent ID 11000033, source code catalogue entry, Question 3.12, PDF page 34. No individual proposer or separate theorem reference is supplied.
\item Farb, \emph{Problems on Mapping Class Groups and Related Topics}. Supplied URL: \url{https://www.math.uchicago.edu/~farb/papers/mcgbook.pdf}; recorded access 2026-07-29; extraction pdf-font-structure. The edited volume contains 329 explicitly typeset problem, conjecture, and question headings across chapters, with repeated local numbering disambiguated by page. The trailing page header ``28 B. Farb'' is extraction context.
\item Supplied status open and OPEN-TRIAGE. The research summary says ``Not settled through 2026; no verification found,'' without a supporting publication or complete audit. Current primary-literature status and redistribution terms for the publicly accessible author/publisher copy remain NEEDS\_REVIEW.
\end{itemize}
\end{document}
